\section{Does the probe track the rule, or the emission?}
\label{app:policy}

The emission.
Changing the operator's rule to permit sharing leaves the probe's score on leaking calls essentially unchanged on both models (\cref{tab:policy}).

\paragraph{Design.}
We take the leaking calls of the fresh test episodes and matched non-leaking calls from the same scenarios and rounds, and re-run each call's exact prompt through the model under three system prompts: the original rule, a rule that explicitly permits sharing the value with teammates, and a neutral rewording that permits nothing new.
The saved reply is held fixed, so only the rule changes, and the frozen last-prompt probe scores the last prompt token as before, with no refitting.
On gpt-oss-20b this is the frozen last-prompt probe, not the pre-registered reasoning probe.

\paragraph{Result.}
Permitting the leak does not detectably lower the score of leaking calls, and the neutral rewording also leaves it unchanged.
The probe ranks leaking above non-leaking calls about equally well under all three prompts, and nearly the same leaking calls cross the frozen threshold.
The probe therefore mostly tracks that the model is about to write the protected value, not that writing it breaks the rule.

\paragraph{Caveat.}
This is an off-policy counterfactual: under the permissive rule the model might not have written exactly the saved reply, and the rule's separate sentence about the consequences of a disclosure was left unchanged.

\begin{table}[h]
  \centering
  \caption{Policy counterfactual on the fresh test episodes. Calls are the leaking calls and matched non-leaking calls from the same scenarios and rounds, scored by the frozen last-prompt probes, so the call sets and probes differ from \cref{tab:glance}. Score changes are medians over leaking calls with 95\% bootstrap intervals.}
  \label{tab:policy}
  \small
  % Generated by generate_tables.py; edit registry table configuration.
\begin{tabular}{p{0.30\linewidth}p{0.30\linewidth}p{0.30\linewidth}}
\toprule
Policy counterfactual & Qwen3-32B & gpt-oss-20b \\
\midrule
Emitting calls & \qwenPolicyLeakingCalls{} & \gptossPolicyLeakingCalls{} \\
Matched nonemitting calls & \qwenPolicyNonleakingCalls{} & \gptossPolicyNonleakingCalls{} \\
Median PERMIT minus FORBID & \qwenPolicyPermitShift{} \qwenPolicyPermitShiftCI{} & \gptossPolicyPermitShift{} \gptossPolicyPermitShiftCI{} \\
Median NEUTRAL minus FORBID & \qwenPolicyNeutralShift{} \qwenPolicyNeutralShiftCI{} & \gptossPolicyNeutralShift{} \gptossPolicyNeutralShiftCI{} \\
FORBID within-group AUROC & \qwenPolicyForbidAUROC{} \qwenPolicyForbidAUROCCI{} & \gptossPolicyForbidAUROC{} \gptossPolicyForbidAUROCCI{} \\
PERMIT within-group AUROC & \qwenPolicyPermitAUROC{} \qwenPolicyPermitAUROCCI{} & \gptossPolicyPermitAUROC{} \gptossPolicyPermitAUROCCI{} \\
NEUTRAL within-group AUROC & \qwenPolicyNeutralAUROC{} \qwenPolicyNeutralAUROCCI{} & \gptossPolicyNeutralAUROC{} \gptossPolicyNeutralAUROCCI{} \\
FORBID threshold crossings & \qwenPolicyForbidCrossingCount{}/\qwenPolicyForbidCrossingTotal{}; \qwenPolicyForbidCrossing{} \qwenPolicyForbidCrossingCI{} & \gptossPolicyForbidCrossingCount{}/\gptossPolicyForbidCrossingTotal{}; \gptossPolicyForbidCrossing{} \gptossPolicyForbidCrossingCI{} \\
PERMIT threshold crossings & \qwenPolicyPermitCrossingCount{}/\qwenPolicyPermitCrossingTotal{}; \qwenPolicyPermitCrossing{} \qwenPolicyPermitCrossingCI{} & \gptossPolicyPermitCrossingCount{}/\gptossPolicyPermitCrossingTotal{}; \gptossPolicyPermitCrossing{} \gptossPolicyPermitCrossingCI{} \\
NEUTRAL threshold crossings & \qwenPolicyNeutralCrossingCount{}/\qwenPolicyNeutralCrossingTotal{}; \qwenPolicyNeutralCrossing{} \qwenPolicyNeutralCrossingCI{} & \gptossPolicyNeutralCrossingCount{}/\gptossPolicyNeutralCrossingTotal{}; \gptossPolicyNeutralCrossing{} \gptossPolicyNeutralCrossingCI{} \\
\bottomrule
\end{tabular}

\end{table}
